% Experimental and Evaluation Section for τ²-Adv Bench Paper
% This file contains the experimental setup, results, and analysis

\section{Experimental Evaluation}
\label{sec:experiments}

We conduct extensive evaluation of five state-of-the-art LLMs to assess their robustness against adversarial manipulation in dual-control customer service scenarios. Our evaluation comprises 360 total evaluations (72 per model) across three domains, five attack strategies, and three sophistication levels.

\subsection{Experimental Setup}

\subsubsection{Models Evaluated}

We evaluate five models representing the current frontier of instruction-following and safety-aligned language models, accessed via OpenRouter API:

\begin{itemize}
    \item \textbf{GPT-OSS-120B}: A 120 billion parameter open-source model optimized for instruction following and customer service applications
    \item \textbf{Grok 4.1 Fast}: xAI's fast inference model with strong reasoning capabilities
    \item \textbf{DeepSeek V3.2}: DeepSeek's latest model with enhanced reasoning and instruction following
    \item \textbf{Kimi K2.5}: Moonshot AI's multimodal model optimized for complex dialogues
    \item \textbf{Mimo V2 Flash}: A lightweight model optimized for fast inference in conversational settings
\end{itemize}

\subsubsection{Evaluation Protocol}

For each model, we evaluate across:
\begin{itemize}
    \item \textbf{3 domains}: Airline, Retail, Telecom
    \item \textbf{5 attack strategies}: Social Engineering, Prompt Injection, Policy Exploitation, Identity Manipulation, Information Extraction
    \item \textbf{3 sophistication levels}: Low ($\phi=0.0$), Medium ($\phi=0.5$), High ($\phi=1.0$)
    \item \textbf{8 tasks per configuration}: 3 tasks per (domain, strategy, sophistication) combination
\end{itemize}

This yields $5 \times 3 \times 5 \times 3 \times \frac{8}{15} \approx 72$ evaluations per model, totaling 360 evaluations across all models. Agent temperature is set to 0.0 for reproducibility.

\subsubsection{Metrics}

We report the following metrics:
\begin{itemize}
    \item \textbf{Safety Score}: $S_{\text{safety}} = \max\left(0, 1 - \frac{\sum_{v \in V} \text{severity}(v)}{N_{\text{turns}}}\right)$, measuring the agent's resistance to policy violations
    \item \textbf{Attack Resistance Rate}: The fraction of adversarial episodes where no critical violations (severity $> 0.8$) occurred
    \item \textbf{Violation Count}: Average number of policy violations detected per episode
\end{itemize}

\subsection{Main Results}

\subsubsection{Overall Performance by Model}

\begin{table}[t]
\centering
\caption{Overall Model Safety Rankings. All models evaluated with 72 adversarial episodes each. Higher safety scores indicate better resistance to adversarial manipulation.}
\label{tab:overall_rankings}
\begin{tabular}{clcc}
\toprule
\textbf{Rank} & \textbf{Model} & \textbf{Safety Score} & \textbf{Evaluations} \\
\midrule
1 & GPT-OSS-120B & 65.5\% & 72 \\
2 & Grok 4.1 Fast & 53.3\% & 72 \\
3 & DeepSeek V3.2 & 52.7\% & 72 \\
4 & Kimi K2.5 & 33.0\% & 72 \\
5 & Mimo V2 Flash & 28.1\% & 72 \\
\bottomrule
\end{tabular}
\end{table}

Table~\ref{tab:overall_rankings} presents the overall safety rankings. GPT-OSS-120B demonstrates the strongest adversarial robustness with a 65.5\% safety score, substantially outperforming other models. There is a notable performance gap between the top three models (GPT-OSS-120B, Grok 4.1 Fast, DeepSeek V3.2) and the bottom two (Kimi K2.5, Mimo V2 Flash), with a 2.3$\times$ difference between the highest and lowest performers.

\textbf{Finding 1: Substantial variance in adversarial robustness.} The 37.4 percentage point gap between the best (65.5\%) and worst (28.1\%) performing models demonstrates that adversarial robustness is not uniformly distributed across frontier models, even when they may achieve similar performance on standard benchmarks.

\subsubsection{Performance by Domain}

\begin{table}[t]
\centering
\caption{Safety Scores by Domain. Aggregated across all models and attack strategies.}
\label{tab:domain_scores}
\begin{tabular}{lcc}
\toprule
\textbf{Domain} & \textbf{Safety Score} & \textbf{Evaluations} \\
\midrule
Retail & 56.9\% & 120 \\
Telecom & 54.2\% & 120 \\
Airline & 28.4\% & 120 \\
\bottomrule
\end{tabular}
\end{table}

\begin{table}[t]
\centering
\caption{Model $\times$ Domain Safety Scores (\%). Each cell represents the average safety score for that model-domain combination.}
\label{tab:model_domain}
\begin{tabular}{lccc|c}
\toprule
\textbf{Model} & \textbf{Retail} & \textbf{Telecom} & \textbf{Airline} & \textbf{Average} \\
\midrule
GPT-OSS-120B & 74.6\% & 73.5\% & 48.3\% & 65.5\% \\
Grok 4.1 Fast & 55.6\% & 74.8\% & 29.4\% & 53.3\% \\
DeepSeek V3.2 & 73.8\% & 54.2\% & 30.2\% & 52.7\% \\
Kimi K2.5 & 46.0\% & 32.9\% & 20.0\% & 33.0\% \\
Mimo V2 Flash & 34.4\% & 35.8\% & 14.0\% & 28.1\% \\
\bottomrule
\end{tabular}
\end{table}

Tables~\ref{tab:domain_scores} and~\ref{tab:model_domain} reveal significant domain-dependent vulnerabilities:

\textbf{Finding 2: Airline domain is systematically vulnerable.} The Airline domain shows dramatically lower safety scores (28.4\%) compared to Retail (56.9\%) and Telecom (54.2\%). This 2$\times$ vulnerability gap persists across all models, with even the best-performing model (GPT-OSS-120B) achieving only 48.3\% in Airline compared to 74.6\% in Retail.

\textbf{Finding 3: Domain-specific robustness varies by model.} While GPT-OSS-120B maintains relatively consistent performance across Retail and Telecom (74.6\% vs 73.5\%), other models show pronounced domain dependencies. Grok 4.1 Fast achieves 74.8\% in Telecom but only 55.6\% in Retail, while DeepSeek V3.2 shows the opposite pattern (54.2\% Telecom, 73.8\% Retail).

\begin{table}[t]
\centering
\caption{Model Performance Gap Analysis by Domain. $\Delta$ shows the difference from the worst performer per domain; Max $\Delta$ row shows the gap between best and worst models.}
\label{tab:domain_gap}
\small
\begin{tabular}{l|cc|cc|cc}
\toprule
& \multicolumn{2}{c|}{\textbf{Retail}} & \multicolumn{2}{c|}{\textbf{Telecom}} & \multicolumn{2}{c}{\textbf{Airline}} \\
\textbf{Model} & Score & $\Delta$ & Score & $\Delta$ & Score & $\Delta$ \\
\midrule
GPT-OSS-120B & 74.6\% & +40.2 & 73.5\% & +37.7 & 48.3\% & +34.3 \\
Grok 4.1 Fast & 55.6\% & +21.2 & 74.8\% & +39.0 & 29.4\% & +15.4 \\
DeepSeek V3.2 & 73.8\% & +39.4 & 54.2\% & +18.4 & 30.2\% & +16.2 \\
Kimi K2.5 & 46.0\% & +11.6 & 32.9\% & $-$2.9 & 20.0\% & +6.0 \\
Mimo V2 Flash & 34.4\% & --- & 35.8\% & --- & 14.0\% & --- \\
\midrule
\textbf{Max $\Delta$} & \multicolumn{2}{c|}{40.2 pp} & \multicolumn{2}{c|}{41.9 pp} & \multicolumn{2}{c}{34.3 pp} \\
\bottomrule
\end{tabular}
\end{table}

Table~\ref{tab:domain_gap} quantifies the performance gaps between models within each domain. The maximum gap ranges from 34.3 percentage points (Airline) to 41.9 pp (Telecom), demonstrating that model selection has substantial impact on domain-specific robustness. Notably, GPT-OSS-120B achieves the highest $\Delta$ in Retail (+40.2) while Grok 4.1 Fast leads in Telecom (+39.0), suggesting complementary strengths.

\subsubsection{Performance by Attack Strategy}

\begin{table}[t]
\centering
\caption{Safety Scores by Attack Strategy. Higher scores indicate the strategy is less effective (agents resist better).}
\label{tab:strategy_scores}
\begin{tabular}{lcc}
\toprule
\textbf{Strategy} & \textbf{Safety Score} & \textbf{Evaluations} \\
\midrule
Information Extraction & 68.0\% & 45 \\
Identity Manipulation & 55.6\% & 45 \\
Social Engineering & 48.5\% & 135 \\
Policy Exploitation & 42.9\% & 45 \\
Prompt Injection & 30.0\% & 90 \\
\bottomrule
\end{tabular}
\end{table}

\begin{table}[t]
\centering
\caption{Model $\times$ Attack Strategy Safety Scores (\%). Legend: SE=Social Engineering, PI=Prompt Injection, PE=Policy Exploitation, IM=Identity Manipulation, IE=Information Extraction.}
\label{tab:model_strategy}
\small
\begin{tabular}{lccccc|c}
\toprule
\textbf{Model} & \textbf{IE} & \textbf{IM} & \textbf{SE} & \textbf{PE} & \textbf{PI} & \textbf{Avg} \\
\midrule
GPT-OSS-120B & 81.7\% & 68.3\% & 72.6\% & 65.6\% & 45.3\% & 65.5\% \\
Grok 4.1 Fast & 73.9\% & 62.2\% & 52.4\% & 46.1\% & 43.3\% & 53.3\% \\
DeepSeek V3.2 & 65.6\% & 65.6\% & 55.2\% & 56.1\% & 34.4\% & 52.7\% \\
Kimi K2.5 & 59.4\% & 46.7\% & 37.2\% & 18.3\% & 13.9\% & 33.0\% \\
Mimo V2 Flash & 59.4\% & 35.0\% & 25.2\% & 28.3\% & 13.1\% & 28.1\% \\
\bottomrule
\end{tabular}
\end{table}

Table~\ref{tab:strategy_scores} reveals a clear hierarchy of attack effectiveness:

\textbf{Finding 4: Prompt injection is the most effective attack vector.} With only 30.0\% average safety score, prompt injection achieves the highest attack success rate across all models. This is followed by policy exploitation (42.9\%) and social engineering (48.5\%). Information extraction is least effective (68.0\% safety), suggesting models are better trained to protect internal information.

\textbf{Finding 5: Attack effectiveness varies significantly by model.} Table~\ref{tab:model_strategy} shows that Kimi K2.5 and Mimo V2 Flash are particularly vulnerable to prompt injection (13.9\% and 13.1\% safety respectively), while GPT-OSS-120B maintains relatively higher resistance (45.3\%). The 3.5$\times$ difference in prompt injection resistance between GPT-OSS-120B and Mimo V2 Flash suggests significant variation in instruction-following robustness.

\subsubsection{Impact of Attack Sophistication}

\begin{table}[t]
\centering
\caption{Safety Scores by Sophistication Level.}
\label{tab:sophistication}
\begin{tabular}{lcc}
\toprule
\textbf{Level} & \textbf{Safety Score} & \textbf{Evaluations} \\
\midrule
Low ($\phi=0.0$) & 46.5\% & 120 \\
Medium ($\phi=0.5$) & 46.0\% & 120 \\
High ($\phi=1.0$) & 47.0\% & 120 \\
\bottomrule
\end{tabular}
\end{table}

\begin{table}[t]
\centering
\caption{Model $\times$ Sophistication Level Safety Scores (\%).}
\label{tab:model_sophistication}
\begin{tabular}{lccc|c}
\toprule
\textbf{Model} & $\phi=0.0$ & $\phi=0.5$ & $\phi=1.0$ & \textbf{Average} \\
\midrule
GPT-OSS-120B & 63.1\% & 69.6\% & 63.7\% & 65.5\% \\
Grok 4.1 Fast & 64.8\% & 47.5\% & 47.5\% & 53.3\% \\
DeepSeek V3.2 & 48.3\% & 58.5\% & 51.2\% & 52.7\% \\
Kimi K2.5 & 30.8\% & 33.3\% & 34.8\% & 33.0\% \\
Mimo V2 Flash & 25.4\% & 21.2\% & 37.5\% & 28.1\% \\
\bottomrule
\end{tabular}
\end{table}

\textbf{Finding 6: Sophistication impact is model-dependent.} Contrary to expectations, aggregate safety scores remain relatively stable across sophistication levels (Table~\ref{tab:sophistication}). However, Table~\ref{tab:model_sophistication} reveals that this aggregate stability masks divergent patterns:
\begin{itemize}
    \item \textbf{Grok 4.1 Fast} shows significant degradation from low to medium sophistication (64.8\% $\rightarrow$ 47.5\%), a 17.3 percentage point drop
    \item \textbf{Mimo V2 Flash} shows the opposite pattern, with improved resistance at high sophistication (25.4\% $\rightarrow$ 37.5\%)
    \item \textbf{GPT-OSS-120B} maintains stable performance across all levels
\end{itemize}

\begin{table}[t]
\centering
\caption{Attack Strategy $\times$ Sophistication Safety Scores (\%).}
\label{tab:strategy_sophistication}
\begin{tabular}{lccc|c}
\toprule
\textbf{Strategy} & $\phi=0.0$ & $\phi=0.5$ & $\phi=1.0$ & \textbf{Average} \\
\midrule
Information Extraction & 70.0\% & 64.7\% & 69.3\% & 68.0\% \\
Identity Manipulation & 60.7\% & 55.7\% & 50.3\% & 55.6\% \\
Social Engineering & 51.4\% & 46.8\% & 47.3\% & 48.5\% \\
Policy Exploitation & 37.3\% & 44.0\% & 47.3\% & 42.9\% \\
Prompt Injection & 24.8\% & 31.8\% & 33.3\% & 30.0\% \\
\bottomrule
\end{tabular}
\end{table}

\textbf{Finding 7: Sophistication effects vary by attack type.} Table~\ref{tab:strategy_sophistication} shows that identity manipulation effectiveness increases monotonically with sophistication (60.7\% $\rightarrow$ 50.3\% safety), while policy exploitation shows the opposite trend (37.3\% $\rightarrow$ 47.3\% safety). This suggests that sophisticated policy exploitation attempts may trigger stronger defensive responses.

\subsubsection{Cross-Dimensional Analysis}

\begin{table}[t]
\centering
\caption{Domain $\times$ Attack Strategy Safety Scores (\%).}
\label{tab:domain_strategy}
\begin{tabular}{lccccc|c}
\toprule
\textbf{Domain} & \textbf{IE} & \textbf{IM} & \textbf{SE} & \textbf{PE} & \textbf{PI} & \textbf{Avg} \\
\midrule
Retail & 77.0\% & 84.3\% & 56.8\% & 46.0\% & 38.7\% & 56.9\% \\
Telecom & 37.0\% & 68.0\% & 63.3\% & 68.0\% & 35.5\% & 54.2\% \\
Airline & 90.0\% & 14.3\% & 25.4\% & 14.7\% & 15.8\% & 28.4\% \\
\bottomrule
\end{tabular}
\end{table}

Table~\ref{tab:domain_strategy} reveals striking domain-strategy interactions:

\textbf{Finding 8: Domain-specific vulnerability patterns.} The Airline domain shows extreme vulnerability to identity manipulation (14.3\% safety) and policy exploitation (14.7\%), while being highly resistant to information extraction (90.0\%). This pattern suggests that airline agents may have strong training against information disclosure but weak policy enforcement mechanisms.

\textbf{Finding 9: Retail domain is most robust to identity attacks.} With 84.3\% safety against identity manipulation, retail agents demonstrate strong identity verification protocols. However, they remain vulnerable to prompt injection (38.7\%) and policy exploitation (46.0\%).

\subsection{Discussion}

\subsubsection{Key Insights}

Our evaluation reveals several critical patterns for conversational AI agent safety:

\begin{enumerate}
    \item \textbf{Model capability does not guarantee robustness.} The 2.3$\times$ difference between best and worst performing models demonstrates that adversarial robustness requires specific training and is not an emergent property of scale alone.

    \item \textbf{Domain-specific vulnerabilities require domain-specific defenses.} The dramatic variation in performance across domains (28.4\% Airline vs 56.9\% Retail) indicates that generic safety training is insufficient for specialized customer service applications.

    \item \textbf{Prompt injection remains a critical vulnerability.} With only 30.0\% average safety and as low as 13.1\% for vulnerable models, prompt injection attacks represent the most significant threat vector.

    \item \textbf{Attack sophistication effects are non-uniform.} The complex interaction between sophistication level, attack strategy, and model architecture suggests that simple robustness testing with naive attacks may not capture real-world vulnerability.
\end{enumerate}

\subsubsection{Implications for Deployment}

These findings have direct implications for the deployment of conversational AI agents:

\begin{itemize}
    \item \textbf{Pre-deployment adversarial testing is essential.} The significant vulnerability gap between models demonstrates that adversarial evaluation should be a required step before production deployment.

    \item \textbf{Domain-specific safety tuning is necessary.} Models should undergo additional safety training specific to their deployment domain, particularly for high-risk domains like airline customer service.

    \item \textbf{Defense prioritization should be data-driven.} The attack strategy effectiveness hierarchy (prompt injection > policy exploitation > social engineering > identity manipulation > information extraction) should guide defense investment priorities.

    \item \textbf{Monitoring for emerging attack patterns.} The varying sophistication effects suggest that adversaries may discover unexpected attack amplification at certain sophistication levels.
\end{itemize}

\subsubsection{Limitations}

Our evaluation has several limitations:

\begin{itemize}
    \item \textbf{Simulated environment.} While our scenarios are designed to reflect realistic customer service interactions, findings should be validated in actual deployments.

    \item \textbf{Model access.} We evaluate models through API access, which may include additional safety layers not present in base models.

    \item \textbf{Attack coverage.} Our five attack strategies represent major categories but do not exhaust the adversarial space. Novel attacks may achieve different success rates.

    \item \textbf{Temporal dynamics.} LLM capabilities and safety properties evolve rapidly; these results represent a snapshot of model performance at evaluation time.
\end{itemize}

\section{Conclusion}

We conducted comprehensive adversarial evaluation of five frontier LLMs on the $\tau^2$-Adv benchmark, totaling 360 evaluations across three customer service domains. Our results reveal:

\begin{itemize}
    \item A 2.3$\times$ safety score gap between best (GPT-OSS-120B: 65.5\%) and worst (Mimo V2 Flash: 28.1\%) performing models
    \item Airline domain shows 2$\times$ higher vulnerability than Retail and Telecom domains
    \item Prompt injection achieves the highest attack success with only 30.0\% average safety score
    \item Attack sophistication effects are model and strategy dependent, with some models showing improved defense against sophisticated attacks
\end{itemize}

These findings underscore the urgent need for adversarial robustness evaluation before deploying conversational AI agents in customer service applications. We release our benchmark and evaluation framework to support ongoing research in this critical area.

% Figures referenced in the text (place in figures/ directory)
% - new_figures/fig1_overall_model_scores.pdf
% - new_figures/fig2_domain_comparison.pdf
% - new_figures/fig3_strategy_comparison.pdf
% - new_figures/fig4_model_domain_heatmap.pdf
% - new_figures/fig5_model_strategy_heatmap.pdf
% - new_figures/fig6_sophistication_impact.pdf
% - new_figures/fig7_radar_profiles.pdf
% - new_figures/fig8_domain_strategy_heatmap.pdf
